\documentclass[11pt,a4paper]{article}

\usepackage[utf8]{inputenc}
\usepackage[T1]{fontenc}
\usepackage{lmodern}
\usepackage[margin=1in]{geometry}
\usepackage{hyperref}
\usepackage{booktabs}
\usepackage{enumitem}
\usepackage{listings}
\usepackage{xcolor}
\usepackage{amsmath}
\usepackage[most]{tcolorbox}

\hypersetup{
  colorlinks=true,
  linkcolor=blue,
  urlcolor=blue,
}

\lstset{
  language=Python,
  basicstyle=\ttfamily\small,
  backgroundcolor=\color{gray!10},
  frame=single,
  rulecolor=\color{gray!30},
  breaklines=true,
  showstringspaces=false,
}

\newtcolorbox{notebox}{
  colback=blue!5,
  colframe=blue!40,
  fonttitle=\bfseries,
  title=Note,
  boxrule=0.5pt,
  arc=2pt,
}

\title{Think Different}
\author{}
\date{}

\begin{document}
\maketitle

\section{The Campaign That Defined Apple}

In 1997, Apple was weeks from bankruptcy. Steve Jobs had just returned, and the company needed to remind the world --- and itself --- what it stood for.

The \emph{Think Different} campaign, created with TBWA\textbackslash{}Chiat\textbackslash{}Day, became one of the most iconic advertising campaigns in history. It didn't feature a single Apple product.

\section{The Misfits and Rebels}

The campaign celebrated people who changed the world by thinking differently:

\begin{itemize}
  \item \textbf{Albert Einstein} --- rewrote the laws of physics
  \item \textbf{Mahatma Gandhi} --- led a nation through nonviolence
  \item \textbf{Martha Graham} --- revolutionized modern dance
  \item \textbf{Pablo Picasso} --- shattered every rule of art
  \item \textbf{Amelia Earhart} --- flew where no one had flown
  \item \textbf{Muhammad Ali} --- floated and stung and never backed down
  \item \textbf{Jim Henson} --- gave felt and foam a soul
  \item \textbf{Maria Callas} --- made opera feel dangerous
  \item \textbf{John Lennon \& Yoko Ono} --- imagined peace into a movement
\end{itemize}

\begin{quote}
The ones who see things differently. They're not fond of rules.
And they have no respect for the status quo.
\end{quote}

\section{The Posters}

The print campaign was striking in its simplicity:

\begin{center}
\begin{tabular}{ll}
  \toprule
  \textbf{Element} & \textbf{Description} \\
  \midrule
  Photo  & Black-and-white portrait \\
  Logo   & Rainbow Apple logo \\
  Text   & Just two words: \emph{Think Different} \\
  Layout & Clean, minimal, iconic \\
  \bottomrule
\end{tabular}
\end{center}

Each poster was a full-page black-and-white photograph with the Apple logo and the words \emph{Think Different} in the corner. Nothing else. No product shots, no specs, no pricing.

\subsection{The Genius of Restraint}

\begin{notebox}
The campaign succeeded \emph{because} of what it left out. In an era of feature-laden tech ads, Apple showed faces instead of products, aspirations instead of specifications.
\end{notebox}

\section{Timeline}

\begin{enumerate}
  \item \textbf{1997} --- Jobs returns to Apple\footnote{Jobs had been ousted from Apple in 1985 and spent 12 years building NeXT and Pixar before returning.}
  \item \textbf{1997} --- Campaign launches with TV spot and print
  \item \textbf{1998} --- iMac launches, riding the wave of renewed brand identity
  \item \textbf{2002} --- Campaign winds down as Apple shifts to product-focused marketing
\end{enumerate}

\section{The Legacy}

The campaign did more than sell computers. It \emph{repositioned Apple's entire brand identity} --- from a struggling tech company to a cultural force.

\subsection{Impact on Design}

Apple's design philosophy --- \texttt{simplicity over complexity} --- traces directly back to \emph{Think Different}. The campaign's aesthetic DNA lives on in:

\begin{itemize}
  \item The clean lines of every Apple Store
  \item The white space of every keynote slide
  \item The minimalism of every product launch
\end{itemize}

\subsection{The Cultural Moment}

\begin{notebox}
Think Different wasn't just advertising. It was a \emph{manifesto}. It told creative people everywhere that a computer company understood them --- and that technology could be a tool for changing the world, not just a beige box on a desk.
\end{notebox}

\section{Technical Notes}

This file demonstrates Djot markup features:

\begin{lstlisting}
# Djot supports fenced code blocks with language tags
def think_different():
    """The ones who are crazy enough to think they can change the world
    are the ones who do."""
    return "Here's to the crazy ones."
\end{lstlisting}

Inline code: \texttt{Think Different} was trademarked in 1997.

Math example: Apple's stock price growth from 1997 to 2002 was approximately $\frac{25 - 3.19}{3.19} \approx 684\%$.

A link to \href{https://apple.com}{Apple's history} and a reference link to \href{https://en.wikipedia.org/wiki/Think_different}{the Wikipedia article}.

\subsection{Footnotes in Action}

The campaign's voiceover was originally recorded by Richard Dreyfuss\footnote{Dreyfuss was chosen for his warm, authoritative tone --- he could make a manifesto sound like a bedtime story.}, though Jobs also recorded a version himself that was never aired publicly.

\bigskip
\noindent\rule{\textwidth}{0.4pt}

\medskip
\noindent\emph{Think Different} --- two words that changed everything.

\end{document}
